\section{Convergence of the hard-margin SVM}
\label{sec:svm}

\subsection{The reduced primal problem}

We consider the hard-margin linear SVM for the rescaled data $\{u_p\}_{p\in\cI}$:
\begin{equation}
\min_{w,b}\ \tfrac12\|w\|^2
\quad\text{subject to}\quad
y_p(w^Tu_p+b)\ge 1\quad\text{for all }p\in\cI.
\label{eq:5.1}
\end{equation}
The dual problem is given by
\begin{equation}
\max_{\alpha\ge 0}\ F(\alpha,G)=\mathbf{1}^T\alpha-\tfrac12\alpha^TYGY\alpha
\quad\text{subject to}\quad y^T\alpha=0,
\label{eq:5.2}
\end{equation}
where $Y=\diag(y)$ and $w=\sum_p\alpha_p y_p u_p$.

In this section, we also use the following reduced form of~\eqref{eq:5.1}, in which
the $d$-dimensional variable $w$ is eliminated. The optimal $w$ belongs to the
span of $\{u_p\}$, because writing $w=w_\parallel+w_\perp$ we have
$\|w\|^2=\|w_\parallel\|^2+\|w_\perp\|^2$ while $w^Tu_p=w_\parallel^Tu_p$. Hence we
may write $w=\sum_p\beta_p u_p$, so that $\|w\|^2=\beta^TG\beta$ and
$w^Tu_p=(G\beta)_p$. Then~\eqref{eq:5.1} is equivalent to
\begin{equation}
\min_{(\beta,b)}\ \tfrac12\beta^TG\beta
\quad\text{subject to}\quad
y_p\bigl((G\beta)_p+b\bigr)\ge 1\quad\text{for all }p\in\cI,
\label{eq:5.3}
\end{equation}
which is a problem in $(\beta,b)\in\R^N\times\R$ and refers to the data only
through $G$. The discriminant function for a new observation $u_0$ is given by
\begin{equation}
f(u_0)=w^Tu_0+b=\sum_{p\in\cI}\beta_p\,u_p^Tu_0+b,
\qquad\beta_p=\alpha_p y_p.
\label{eq:5.4}
\end{equation}
We note that $\beta=Y\alpha$.

\begin{lemma}
\label{lem:4}
Assume $G\succ O$. Then $\{u_p\}$ is affinely independent. Moreover, for any
labelling, \eqref{eq:5.1} and~\eqref{eq:5.3} are feasible and their solutions are
unique, and the solution of~\eqref{eq:5.2} is unique.
\end{lemma}

\begin{proof}
The condition $G\succ O$ is equivalent to the linear independence of $\{u_p\}$,
because $a^TGa=\|\sum a_p u_p\|^2$ implies $\sum a_p u_p=0$ and hence $a=0$. The
linear independence implies the affine independence. For affinely independent $N$
points, there exists $(w,b)$ such that $y_p(w^Tu_p+b)=1$ for all $p$, so that the
problem is feasible. The objective function $\tfrac12\beta^TG\beta$ is strictly
convex in $\beta$ from $G\succ O$ and the constraints are convex, and $b$ is
determined uniquely from the active constraints once $\beta$ is given. Hence the
solution of~\eqref{eq:5.3} is unique. The dual objective function is strictly
concave from $YGY\succ O$, so that the solution of~\eqref{eq:5.2} is unique.
\end{proof}

\subsection{Continuity of the solution map}

In Lemma~\ref{lem:5} below, we investigate how the solution of the optimization
problem moves when the parameter $G$ moves. We cannot write the solution
explicitly, because the set of the active constraints changes with $G$. On the
other hand, what we need is not the differentiability but only the continuity. For
this purpose we use the maximum theorem of \citet{berge1963}, which we summarize
in Appendix~\ref{app:A.1} together with its intuitive meaning. Roughly speaking,
the theorem states that the optimal value moves continuously while the optimal
solution may jump, but the point to which it jumps always belongs to the solution
set of the limiting problem. When the solution is unique, this gives the continuity
of the solution map.

\begin{lemma}
\label{lem:5}
Let $\cG_+=\{G\in\R^{N\times N}:G=G^T,\ G\succ O\}$. Then the maps
$G\mapsto\alpha(G)$ and $G\mapsto(\beta(G),b(G))$ are continuous on $\cG_+$,
where $\alpha(G)$ denotes the solution of~\eqref{eq:5.2} and $(\beta(G),b(G))$
denotes the solution of~\eqref{eq:5.3}.
\end{lemma}

\begin{proof}
\textit{Step~1.} Let $G_0\in\cG_+$ and $\eta=\lambda_{\min}(G_0)/2$. Since the
eigenvalues are continuous functions of the matrix, there exists an open
neighbourhood $U$ of $G_0$ such that $\lambda_{\min}(G)\ge\eta$ for all $G\in U$.

\textit{Step~2.} We give an a priori bound of the dual variable. Since $Y$ is an
orthogonal matrix, we have $\alpha^TYGY\alpha=\|G^{1/2}Y\alpha\|^2$, so that
\[
\alpha^TYGY\alpha\ge\eta\|\alpha\|^2
\quad\text{for all }G\in U.
\]
Together with $\mathbf{1}^T\alpha\le\sqrt{N}\,\|\alpha\|$, we have
\begin{equation}
F(\alpha,G)\le\sqrt{N}\,\|\alpha\|-\frac{\eta}{2}\|\alpha\|^2.
\label{eq:5.5}
\end{equation}
Since $\alpha=0$ is feasible and $F(0,G)=0$, the optimal solution satisfies
$F(\alpha,G)\ge 0$. From~\eqref{eq:5.5} we obtain
$\|\alpha\|\le 2\sqrt{N}/\eta$, that is,
\begin{equation}
\|\alpha\|\le\frac{2\sqrt{N}}{\eta}=:K_0
\quad\text{for all }G\in U.
\label{eq:5.6}
\end{equation}
We emphasize that $K_0$ does not depend on $G\in U$.

\textit{Step~3.} From~\eqref{eq:5.6}, we may restrict the feasible set of
\eqref{eq:5.2} to
\[
\cA=\{\alpha\ge 0:\ y^T\alpha=0,\ \|\alpha\|\le K_0\}
\]
without changing the optimal solution on $U$. The set $\cA$ is compact and does
not depend on $G$, that is, the feasible correspondence is a constant
correspondence, which is trivially continuous. On the other hand, $F(\alpha,G)$
is a polynomial in $(\alpha,G)$ and hence continuous. Thus all the conditions of
the maximum theorem are met, and the solution correspondence is upper
hemicontinuous on $U$. Since the solution is a singleton from Lemma~\ref{lem:4},
the solution map is continuous on $U$ from Corollary~\ref{cor:A.1}. Since $G_0$
is arbitrary, the map is continuous on $\cG_+$.

\textit{Step~4.} We consider~\eqref{eq:5.3}. The optimal value is bounded above by
the value at the feasible point given in the proof of Lemma~\ref{lem:4}, which is
a continuous function of $G$ and hence bounded on $U$, say by $M_U$. Then
\[
\tfrac12\beta^TG\beta\le M_U
\quad\Rightarrow\quad
\|\beta\|^2\le\frac{2M_U}{\eta}.
\]
Moreover, at the optimal solution, at least one constraint is active, since
otherwise one could shrink $(\beta,b)$ and decrease the objective function. For
such $p$ we have $|(G\beta)_p+b|=1$, so that
$|b|\le 1+\|G\|\|\beta\|$. Hence we may restrict the feasible set to a compact set
which does not depend on $G\in U$, and the same argument as in Step~3 gives the
continuity.
\end{proof}

\subsection{Convergence of the solution}

\begin{theorem}
\label{thm:2}
Assume \textup{(A1)} to \textup{(A4)} and \textup{(S)}. Then, as $d\to\infty$,
it holds that
\[
\alpha^{(d)}\wto\alpha^*=\alpha(G^*),
\qquad
(\beta^{(d)},b^{(d)})\wto(\beta^*,b^*)=(\beta(G^*),b(G^*)),
\]
where $G^*$ is the limiting Gram matrix given in Theorem~\ref{thm:1}.
\end{theorem}

\begin{proof}
From Theorem~\ref{thm:1} we have $G\wto G^*$ and $P(G^*\succ O)=1$. From
Lemma~\ref{lem:5}, the maps are continuous on $\cG_+$, so that the set of the
discontinuity points has measure zero under the limiting distribution. Thus we
obtain the result from the continuous mapping theorem.
\end{proof}
